\section{Introduction}
\label{sec:intro}

\IEEEPARstart{R}{iver} and reservoir monitoring with in-situ sensor nodes has to work where field visits are expensive and power is scarce~\cite{adumanu2017survey,adumanu2020smart}. A typical node samples temperature (T), specific conductance (EC), dissolved oxygen (DO) and pH every 5--30~minutes and sends the readings over a cellular or low-power wide-area network (LPWAN). Such nodes lose data. The InWaterSense deployment on the Sitnica river in Kosovo delivered \inwsDelivered{} of its expected readings and missed \inwsMissing{} (\inwsLoss)~\cite{ahmedi2022dataset}; the CSIRO water-quality network in Australia averaged about 64\% delivery because the gateway's connection to the server ``was down occasionally''~\cite{ledinh2007design}; and deployment reports describe gateways that lose power every night~\cite{langendoen2006murphy}, sinks that stop reconnecting until someone reboots them~\cite{ingelrest2010sensorscope,brzozowski2018lessons}, outages lasting days to weeks~\cite{puluckul2025subarctic}, and prolonged outages of cellular backhaul in remote areas~\cite{zennaro2025kyrgyzstan}.

There are three standard responses to a lost reading. A node can \emph{retransmit} it until it is acknowledged (automatic repeat request, ARQ); it can \emph{buffer} it and \emph{backfill} it later; or the server can \emph{reconstruct} it from the node's own history or from neighbouring sites. Which response is right depends on how data are actually lost. Retransmission helps against a short radio fade but spends energy in vain against a gateway that is down for a day; reconstruction can fill a ten-minute gap but not a week. Yet the analyses behind these remedies assume losses that are short: estimation over lossy links is analysed with independent or short-memory (Markov fading) packet errors~\cite{sinopoli2004kalman,huang2019retransmit,huang2020harq}, and state-of-the-art imputation methods are evaluated on random masks or on simulated failures of a few hours~\cite{kong2014dataloss,cini2022grin,marisca2022spin,chang2023deepcs}. Kong et al.~\cite{kong2014dataloss} describe long losses caused by damaged or depleted nodes, but as named patterns, not as measured durations. How long real delivery gaps last, how much of the missing data they hold, and what that implies for the energy and timeliness of each remedy have not been measured for environmental telemetry.

The contribution of this paper is not a new retransmission or prediction primitive: ARQ, buffering, back-off, batching and error-bounded (dual-prediction) reporting are all established. It is an evidence-to-design chain for outage-aware, energy-efficient telemetry. Real delivery records are used to measure how long losses last and where the missing data are. The records are then replayed through a measurement-parameterised LPWAN energy model, with explicit bounds on what the records cannot tell (why a reading is missing, and when within a slot the network returned), to decide when retransmission, buffering, batching and reconstruction pay. Finally, the reporting design that follows from these findings is fixed on one year of river data and tested on a held-out year. Four research questions structure the paper:
\begin{enumerate}
\item[RQ1] In the records studied, how long are the gaps in delivery, how much of the missing data do long gaps hold, and is a geometric (memoryless) run-length model adequate?
\item[RQ2] Under these gaps, what do retransmission, buffering with back-off and batching each recover, and at what energy and delay on LoRaWAN and NB-IoT?
\item[RQ3] Can reconstruction, from a site's own history or from neighbouring sites on the same river, replace delivery during a gap?
\item[RQ4] Can one reporting scheme raise alarms as fast as per-reading reporting and keep the server's real-time picture within a stated tolerance as often, while using substantially less energy than per-reading retransmission, on data not used to design it?
\end{enumerate}

The contributions follow the chain.
\begin{itemize}
\item \emph{Empirical loss structure} (Sec.~\ref{sec:outages}). In four records, two water-quality deployments and two LoRaWAN networks, losses separate into short radio losses and rare long delivery gaps or system outages, and in the water-quality records gaps of 6~h or more hold most of the missing readings. A geometric run-length model is strongly unsupported relative to the alternatives tested, and fails worst in the long-gap tail that decides between remedies. Where the records allow it, the cause of the long gaps is examined: node-wide silence on the single-node record, and mostly a shared gateway or backhaul path on the two-node record.
\item \emph{Consequence analysis with explicit bounds} (Secs.~\ref{sec:handling} and \ref{sec:reconstruct}). Replaying the records shows that ARQ recovers under a fifth of the missing readings even under optimistic assumptions on retry timing, at \arqRatioMin--\arqRatioMax{} the energy of unconfirmed reporting on LoRaWAN, whereas buffering with probe back-off recovers every reading the node generated at an energy close to that of recovering nothing in the model, and within \mcBufLo--\mcBufHi{} of it under wide parameter uncertainty. The results are bracketed for the unknown cause of gaps, the timing of retries within a slot, finite storage and the duty cycle, the back-off cap, and many nodes recovering at once. For reconstruction, the evaluated estimators, from a site's own day--night (diel) pattern to two contemporary deep imputers, keep few day-long gaps within tolerance, and a test against the last value alone overstates the value of neighbouring sites.
\item \emph{An outage-aware reporting design, confirmed on held-out data} (Sec.~\ref{sec:design}). Bounded event-driven batching combines hourly batches, immediate uploads on a node-side alarm, a bounded-error trigger that mirrors the server's estimate, and buffering with back-off. It was developed on \designSites{} USGS river sites in 2023 and then evaluated, with its parameters and decision rules frozen, on \confSites{} sites in 2024.
\end{itemize}

\emph{Scope.} The loss evidence comes from four records, and the two water-quality records cannot fully separate a network outage from a node that stopped producing readings; the paper therefore speaks of delivery gaps, and treats buffering's recovery as an upper bound for gaps of unknown cause. Energy comes from a model parameterised by published measurements, not from hardware built for this study; the claims are ratios between strategies on the same radio, tested against wide parameter uncertainty. The design is evaluated per node, with a separate multi-node model of recovery after a shared outage.

The rest of the paper is organised as follows. Sec.~\ref{sec:related} reviews related work, and Sec.~\ref{sec:data} describes the data, models and assumptions. Secs.~\ref{sec:outages}--\ref{sec:design} answer RQ1--RQ4 in turn, and Sec.~\ref{sec:discussion} discusses practical guidance, validity and limitations.
